% Options for packages loaded elsewhere
\PassOptionsToPackage{unicode}{hyperref}
\PassOptionsToPackage{hyphens}{url}
\documentclass[
  12pt,
  a4paper,
]{article}
\usepackage{xcolor}
\usepackage[margin=1in]{geometry}
\usepackage{amsmath,amssymb}
\setcounter{secnumdepth}{-\maxdimen} % remove section numbering
\usepackage{iftex}
\ifPDFTeX
  \usepackage[T1]{fontenc}
  \usepackage[utf8]{inputenc}
  \usepackage{textcomp} % provide euro and other symbols
\else % if luatex or xetex
  \usepackage{unicode-math} % this also loads fontspec
  \defaultfontfeatures{Scale=MatchLowercase}
  \defaultfontfeatures[\rmfamily]{Ligatures=TeX,Scale=1}
\fi
\usepackage{lmodern}
\ifPDFTeX\else
  % xetex/luatex font selection
\fi
% Use upquote if available, for straight quotes in verbatim environments
\IfFileExists{upquote.sty}{\usepackage{upquote}}{}
\IfFileExists{microtype.sty}{% use microtype if available
  \usepackage[]{microtype}
  \UseMicrotypeSet[protrusion]{basicmath} % disable protrusion for tt fonts
}{}
\makeatletter
\@ifundefined{KOMAClassName}{% if non-KOMA class
  \IfFileExists{parskip.sty}{%
    \usepackage{parskip}
  }{% else
    \setlength{\parindent}{0pt}
    \setlength{\parskip}{6pt plus 2pt minus 1pt}}
}{% if KOMA class
  \KOMAoptions{parskip=half}}
\makeatother
\usepackage{longtable,booktabs,array}
\usepackage{calc} % for calculating minipage widths
% Correct order of tables after \paragraph or \subparagraph
\usepackage{etoolbox}
\makeatletter
\patchcmd\longtable{\par}{\if@noskipsec\mbox{}\fi\par}{}{}
\makeatother
% Allow footnotes in longtable head/foot
\IfFileExists{footnotehyper.sty}{\usepackage{footnotehyper}}{\usepackage{footnote}}
\makesavenoteenv{longtable}
% definitions for citeproc citations
\NewDocumentCommand\citeproctext{}{}
\NewDocumentCommand\citeproc{mm}{%
  \begingroup\def\citeproctext{#2}\cite{#1}\endgroup}
\makeatletter
 % allow citations to break across lines
 \let\@cite@ofmt\@firstofone
 % avoid brackets around text for \cite:
 \def\@biblabel#1{}
 \def\@cite#1#2{{#1\if@tempswa , #2\fi}}
\makeatother
\newlength{\cslhangindent}
\setlength{\cslhangindent}{1.5em}
\newlength{\csllabelwidth}
\setlength{\csllabelwidth}{3em}
\newenvironment{CSLReferences}[2] % #1 hanging-indent, #2 entry-spacing
 {\begin{list}{}{%
  \setlength{\itemindent}{0pt}
  \setlength{\leftmargin}{0pt}
  \setlength{\parsep}{0pt}
  % turn on hanging indent if param 1 is 1
  \ifodd #1
   \setlength{\leftmargin}{\cslhangindent}
   \setlength{\itemindent}{-1\cslhangindent}
  \fi
  % set entry spacing
  \setlength{\itemsep}{#2\baselineskip}}}
 {\end{list}}
\usepackage{calc}
\newcommand{\CSLBlock}[1]{\hfill\break\parbox[t]{\linewidth}{\strut\ignorespaces#1\strut}}
\newcommand{\CSLLeftMargin}[1]{\parbox[t]{\csllabelwidth}{\strut#1\strut}}
\newcommand{\CSLRightInline}[1]{\parbox[t]{\linewidth - \csllabelwidth}{\strut#1\strut}}
\newcommand{\CSLIndent}[1]{\hspace{\cslhangindent}#1}
\setlength{\emergencystretch}{3em} % prevent overfull lines
\providecommand{\tightlist}{%
  \setlength{\itemsep}{0pt}\setlength{\parskip}{0pt}}
\usepackage{bookmark}
\IfFileExists{xurl.sty}{\usepackage{xurl}}{} % add URL line breaks if available
\urlstyle{same}
\hypersetup{
  hidelinks,
  pdfcreator={LaTeX via pandoc}}

\author{}
\date{}

\begin{document}

\section{Digital twins for software engineers: a survey through a potted
plant}\label{digital-twins-for-software-engineers-a-survey-through-a-potted-plant}

\subsection{Scope}\label{scope}

This survey maps the digital twin (DT) literature for one reader: a
software engineer at home with message brokers, containers and
continuous integration (CI), who has been handed a digital twin to build
or assess. It covers what separates a twin from the systems you already
ship, the architectural and deployment choices on offer, the platforms
that claim to spare you the work, the failure modes that recur, and the
exemplars worth cloning. It excludes the physics inside the models,
manufacturing process detail, control theory, and adoption economics.
One example anchors every section: a potted plant, a soil-moisture
probe, and a small pump.

\subsection{1. The connection is the twin, not the
model}\label{the-connection-is-the-twin-not-the-model}

The literature's sharpest distinction is about coupling, not fidelity.
Kritzinger et al.~grade digital counterparts by level of integration,
separating a Digital Model from a Digital Shadow and from a Digital
Twin, and find that work at the highest stage is the scarcest {[}1{]}.
Ellwein et al.~restate the criterion as the degree of \emph{automated}
integration: a Digital Model needs a person to carry changes across by
hand {[}2{]}. A Digital Shadow has ``an automated one-way data flow''
from the physical object to the digital one {[}3{]}, and shadow and twin
are treated as a paired distinction in manufacturing {[}4{]}.

Apply that to the pot. A service that ingests moisture readings and
draws a chart is a digital shadow. It becomes a twin when the same
service decides to run the pump, and the soil gets wetter because the
software said so.

Do not expect the term to carry that meaning on someone else's
requirements document. Definitions vary widely enough that the variance
is itself a reported problem {[}5{]}, {[}6{]}, and a twin reflects only
the properties that matter ``within a specific application context''
{[}7{]} --- which is what licenses you to model a living plant as three
numbers. The sources also disagree on the central point: the Digital
Twin Consortium definition quoted by Bhandal et al.~asks for a virtual
representation ``synchronized at a specified frequency and fidelity''
and names no automated flow back to the physical side {[}8{]}. Under
Kritzinger, your read-only moisture logger is not a twin; under the
Consortium, it is.

\subsection{2. Architecture: familiar shapes, one unfamiliar
edge}\label{architecture-familiar-shapes-one-unfamiliar-edge}

The software architecture community has claimed twins as its own
subject, with a growing body of peer-reviewed architectural solutions
{[}9{]}. Two pattern catalogs offer shapes to instantiate rather than
principles to derive: one for twin architecture {[}10{]}, one for giving
twin models behaviour rather than structure alone {[}11{]}.

The live architectural split is whether a twin is data or a process.
Commercial platforms represent twins as passive JSON entities, while
research increasingly treats them as orchestrated microservices {[}12{]}
--- realised through microservices and serverless functions across the
cloud-to-edge continuum {[}13{]}, or on Kubernetes with ordinary
cloud-native tooling {[}14{]}. Deployments span centralised cloud
platforms, edge-deployable brokers and containerised twins, and managing
them across that range remains under-explored {[}15{]}.

The unfamiliar edge is that one of your services now writes to a pump.
That makes the twin a security boundary: an adversary can attack from
the digital side to reach physical assets, and the attack surface is
large because the paradigm joins two worlds {[}16{]}. Near-real-time
twin traffic can carry critical signals, and adding a twin expands an
already growing attack surface {[}17{]}.

\subsection{3. Buy, compose, or write it}\label{buy-compose-or-write-it}

Reuse is the field's answer to twin construction cost. Talasila et
al.~argue that building a twin is hard because of the variety of assets,
models, data and services to marshal, and propose a generic platform
assembling twins from reusable components offered as a service {[}18{]},
later demonstrating composition of twins from parts {[}19{]}. Gil et
al.~survey the open-source frameworks through a case study {[}20{]}, and
this Digital-Twin-as-a-Service (DTaaS) framing recurs independently
elsewhere {[}21{]}, {[}22{]}.

That last survey supplies vocabulary worth adopting: \emph{twin
fidelity} is how completely the twin represents its counterpart's state
and structure, \emph{twin granularity} is which parts get modelled at
all {[}22{]}. ``Should the plant model include root depth'' is a
granularity question; ``how close should the predicted moisture curve
sit to the probe'' is a fidelity one. Both are easier to answer than
``how good should the model be''.

\subsection{4. Where these projects actually
fail}\label{where-these-projects-actually-fail}

Three failure modes recur, and each is a familiar distributed-systems
problem wearing new clothes.

\textbf{Clocks.} Coupling a simulation to hardware forces simulation
time to follow wall-clock time, making the execution time of a single
simulation iteration a critical parameter {[}23{]}. A plant model that
takes 90 seconds to compute a 60-second step will never close the loop.

\textbf{Update cadence.} Zhang et al.~put both sides of the cost
plainly: too frequent updates burn communication and computation, while
delayed updates let the model drift from reality {[}24{]}. Heithoff et
al. add that consuming everything the sensors produce would stop the
software delivering results in meaningful time {[}3{]}. Your probe's
sampling rate is a design decision, not a device property, and the link
carrying the samples is itself engineered {[}25{]}.

\textbf{Trust.} The simulation-reality gap is named as the main
challenge across worked case studies {[}26{]}, and organisations often
cannot tell how well a twin matches reality or whether to rely on it for
critical decisions {[}27{]}. If the twin waters the plant while you are
away, that judgement is yours.

\subsection{5. Twins inside the development
pipeline}\label{twins-inside-the-development-pipeline}

A twin is a long-lived deployed system, and the literature treats its
operations accordingly. TwinOps connects twins to DevOps practice
{[}28{]}; Barbie tests twin prototypes through automated integration
tests in a CI pipeline, communication protocols included, with a
smart-farming case study and open artefacts for replication {[}29{]};
Dalibor et al.~map the software engineering literature on twins across
domains {[}30{]}. Twins also change under new requirements, which
motivates notation for their continuous evolution {[}31{]} and
automation of their lifecycle management {[}32{]}. Kamburjan et
al.~frame twin and counterpart together as a self-adaptive system: a
greenhouse plant that becomes infected needs a different watering
regime, so the physical system's stage changes which twin components
apply {[}33{]}.

\subsection{6. Comparison: where to
start}\label{comparison-where-to-start}

\begin{longtable}[]{@{}
  >{\raggedright\arraybackslash}p{(\linewidth - 6\tabcolsep) * \real{0.2500}}
  >{\raggedright\arraybackslash}p{(\linewidth - 6\tabcolsep) * \real{0.2500}}
  >{\raggedright\arraybackslash}p{(\linewidth - 6\tabcolsep) * \real{0.2500}}
  >{\raggedright\arraybackslash}p{(\linewidth - 6\tabcolsep) * \real{0.2500}}@{}}
\toprule\noalign{}
\begin{minipage}[b]{\linewidth}\raggedright
Starting point
\end{minipage} & \begin{minipage}[b]{\linewidth}\raggedright
Citekey
\end{minipage} & \begin{minipage}[b]{\linewidth}\raggedright
Core idea
\end{minipage} & \begin{minipage}[b]{\linewidth}\raggedright
Stated limitation
\end{minipage} \\
\midrule\noalign{}
\endhead
\bottomrule\noalign{}
\endlastfoot
Architecture pattern catalog & \texttt{tekinerdogan\_systems\_2020} & A
catalog of twin architecture patterns to instantiate & Patterns for
structure; behaviour is addressed separately \\
Behaviour pattern catalog & \texttt{lehner\_pattern\_2023} & Augments
twin models with behaviour via model-driven engineering & Assumes a
modelling-layer/realisation-layer split \\
Entanglement-aware middleware &
\texttt{bellavista\_entanglement-aware\_2024} & Middleware that keeps
twin and asset coupled, rather than storing twin state & Contrasts
itself with commercial platforms it does not replace \\
Serverless/Kubernetes platform & \texttt{wermann\_ktwin\_2024} &
Cloud-native tooling, open standards for models & Prototype-stage
evaluation \\
Open-source framework survey & \texttt{gil\_survey\_2024} & Compares
frameworks through one case study & Single case limits generalisation \\
DTaaS platform & \texttt{talasila\_realising\_2024},
\texttt{talasila\_composable\_2025} & One tenant-facing platform hosting
many twins built from shared parts & Reuse depends on components
existing for your assets \\
Plant-and-pump exemplar & \texttt{kamburjan\_greenhousedt\_2024} &
Plants, sensors and water pumps as the physical system; asset model in a
knowledge base & A deliberately simple, low-cost greenhouse \\
Tutorial-grade exemplar & \texttt{gomes\_digital\_2025} & Incubator twin
using ordinary differential equations, Kalman filtering and a
service-oriented architecture & Single case study \\
Closed-loop exemplar & \texttt{goffi\_engineering\_2025} & Direct
control of a physical process as a finite state machine &
Domain-specific process constraints \\
\end{longtable}

For this use case, GreenhouseDT is the closest published match: an
extensible architecture whose physical system is plants, sensors and
water pumps, with the asset model held in the twin's knowledge base
{[}34{]}.

\subsection{7. Gaps in this corpus}\label{gaps-in-this-corpus}

Three gaps stand out. First, scale: after reformulating queries across
agriculture, greenhouse, sensor-actuator and small-exemplar phrasings,
only a handful of sources sit at consumer scale {[}19{]},
{[}34{]}--{[}36{]}. Manufacturing and infrastructure dominate, so
guidance about one pot is extrapolated from factories. Second, the
disagreement in §1 is unresolved in the literature, not merely unsettled
here: whether an automated write path to the physical side is required
separates Kritzinger's grading {[}1{]}, {[}2{]} from the Digital Twin
Consortium's synchronisation-only definition {[}8{]}, and no retrieved
source reconciles them. Third, no retrieved source addresses the running
cost --- compute, energy, maintenance --- of keeping a twin synchronised
over years, which for a single plant may exceed the value of the
decisions it makes.

\subsection{References}\label{references}

\protect\phantomsection\label{refs}
\begin{CSLReferences}{0}{0}
\bibitem[\citeproctext]{ref-kritzinger_digital_2018}
\CSLLeftMargin{{[}1{]} }%
\CSLRightInline{W. Kritzinger, M. Karner, G. Traar, J. Henjes, and W.
Sihn, {``Digital {Twin} in manufacturing: {A} categorical literature
review and classification,''} \emph{Ifac-PapersOnline}, vol. 51, no. 11,
pp. 1016--1022, 2018, Accessed: Dec. 11, 2024. {[}Online{]}. Available:
\url{https://www.sciencedirect.com/science/article/pii/S2405896318316021}}

\bibitem[\citeproctext]{ref-ellwein_rethinking_2025}
\CSLLeftMargin{{[}2{]} }%
\CSLRightInline{C. Ellwein \emph{et al.}, {``Rethinking {Asset}
{Administration} {Shell} {Communication} {Types}: {A} {Systematic}
{Mapping} {Study} and {Portfolio}-{Based} {Classification},''}
\emph{Production Engineering}, vol. 20, Dec. 2025, doi:
\href{https://doi.org/10.1007/s11740-025-01378-3}{10.1007/s11740-025-01378-3}.}

\bibitem[\citeproctext]{ref-heithoff_model-based_2024}
\CSLLeftMargin{{[}3{]} }%
\CSLRightInline{M. Heithoff, N. Jansen, J. Michael, F. Rademacher, and
B. Rumpe, {``Model-{Based} {Engineering} of {Multi}-{Purpose} {Digital}
{Twins} in {Manufacturing},''} in \emph{Digital {Twin}: {Fundamentals}
and {Applications}}, S. Sabri, K. Alexandridis, and N. Lee, Eds., Cham:
Springer Nature Switzerland, 2024, pp. 89--126. doi:
\href{https://doi.org/10.1007/978-3-031-67778-6_5}{10.1007/978-3-031-67778-6\_5}.}

\bibitem[\citeproctext]{ref-bergs_concept_2021}
\CSLLeftMargin{{[}4{]} }%
\CSLRightInline{T. Bergs, S. Gierlings, T. Auerbach, A. Klink, D.
Schraknepper, and T. Augspurger, {``The concept of digital twin and
digital shadow in manufacturing,''} \emph{Procedia CIRP}, vol. 101, pp.
81--84, 2021, Accessed: Jan. 21, 2025. {[}Online{]}. Available:
\url{https://www.sciencedirect.com/science/article/pii/S2212827121006612}}

\bibitem[\citeproctext]{ref-vanderhorn_digital_2021}
\CSLLeftMargin{{[}5{]} }%
\CSLRightInline{E. VanDerHorn and S. Mahadevan, {``Digital {Twin}:
{Generalization}, characterization and implementation,''} \emph{Decision
Support Systems}, vol. 145, p. 113524, Jun. 2021, doi:
\href{https://doi.org/10.1016/j.dss.2021.113524}{10.1016/j.dss.2021.113524}.}

\bibitem[\citeproctext]{ref-semeraro_digital_2021}
\CSLLeftMargin{{[}6{]} }%
\CSLRightInline{C. Semeraro, M. Lezoche, H. Panetto, and M. Dassisti,
{``Digital twin paradigm: {A} systematic literature review,''}
\emph{Computers in Industry}, vol. 130, p. 103469, Sep. 2021, doi:
\href{https://doi.org/10.1016/j.compind.2021.103469}{10.1016/j.compind.2021.103469}.}

\bibitem[\citeproctext]{ref-minerva_digital_2020}
\CSLLeftMargin{{[}7{]} }%
\CSLRightInline{R. Minerva, G. M. Lee, and N. Crespi, {``Digital twin in
the {IoT} context: {A} survey on technical features, scenarios, and
architectural models,''} \emph{Proceedings of the IEEE}, vol. 108, no.
10, pp. 1785--1824, 2020, Accessed: Feb. 25, 2025. {[}Online{]}.
Available: \url{https://ieeexplore.ieee.org/abstract/document/9120192/}}

\bibitem[\citeproctext]{ref-bhandal_conceptualising_2024}
\CSLLeftMargin{{[}8{]} }%
\CSLRightInline{R. Bhandal, {``Conceptualising the {Application} of
{Digital} {Twins} in {Supply} {Chain} {Management}: {A} {Path} {Towards}
{Supply} {Chain} {Resilience},''} in \emph{Digital {Twin}:
{Fundamentals} and {Applications}}, S. Sabri, K. Alexandridis, and N.
Lee, Eds., Cham: Springer Nature Switzerland, 2024, pp. 173--189. doi:
\href{https://doi.org/10.1007/978-3-031-67778-6_8}{10.1007/978-3-031-67778-6\_8}.}

\bibitem[\citeproctext]{ref-ferko_architecting_2022}
\CSLLeftMargin{{[}9{]} }%
\CSLRightInline{E. Ferko, A. Bucaioni, and M. Behnam, {``Architecting
{Digital} {Twins},''} \emph{IEEE Access}, vol. 10, pp. 50335--50350,
2022, doi:
\href{https://doi.org/10.1109/ACCESS.2022.3172964}{10.1109/ACCESS.2022.3172964}.}

\bibitem[\citeproctext]{ref-tekinerdogan_systems_2020}
\CSLLeftMargin{{[}10{]} }%
\CSLRightInline{B. Tekinerdogan and C. Verdouw, {``Systems
{Architecture} {Design} {Pattern} {Catalog} for {Developing} {Digital}
{Twins},''} \emph{Sensors}, vol. 20, no. 18, p. 5103, Jan. 2020, doi:
\href{https://doi.org/10.3390/s20185103}{10.3390/s20185103}.}

\bibitem[\citeproctext]{ref-lehner_pattern_2023}
\CSLLeftMargin{{[}11{]} }%
\CSLRightInline{D. Lehner, S. Sint, M. Eisenberg, and M. Wimmer, {``A
pattern catalog for augmenting {Digital} {Twin} models with behavior,''}
\emph{at - Automatisierungstechnik}, vol. 71, no. 6, pp. 423--443, Jun.
2023, doi:
\href{https://doi.org/10.1515/auto-2022-0144}{10.1515/auto-2022-0144}.}

\bibitem[\citeproctext]{ref-bellavista_entanglement-aware_2024}
\CSLLeftMargin{{[}12{]} }%
\CSLRightInline{P. Bellavista, N. Bicocchi, M. Fogli, C. Giannelli, M.
Mamei, and M. Picone, {``An {Entanglement}-{Aware} {Middleware} for
{Digital} {Twins},''} \emph{ACM Trans. Internet Things}, Oct. 2024, doi:
\href{https://doi.org/10.1145/3699520}{10.1145/3699520}.}

\bibitem[\citeproctext]{ref-bellavista_exploiting_2024}
\CSLLeftMargin{{[}13{]} }%
\CSLRightInline{P. Bellavista, N. Bicocchi, M. Fogli, C. Giannelli, M.
Mamei, and M. Picone, {``Exploiting microservices and serverless for
{Digital} {Twins} in the cloud-to-edge continuum,''} \emph{Future
Generation Computer Systems}, vol. 157, pp. 275--287, Aug. 2024, doi:
\href{https://doi.org/10.1016/j.future.2024.03.052}{10.1016/j.future.2024.03.052}.}

\bibitem[\citeproctext]{ref-wermann_ktwin_2024}
\CSLLeftMargin{{[}14{]} }%
\CSLRightInline{A. G. Wermann and J. A. Wickboldt, {``{KTWIN}: {A}
{Serverless} {Kubernetes}-based {Digital} {Twin} {Platform}.''} arXiv,
Aug. 2024. doi:
\href{https://doi.org/10.48550/arXiv.2408.01635}{10.48550/arXiv.2408.01635}.}

\bibitem[\citeproctext]{ref-barbone_digital_2024}
\CSLLeftMargin{{[}15{]} }%
\CSLRightInline{A. Barbone, S. Burattini, M. Martinelli, M. Picone, A.
Ricci, and A. Virdis, {``Digital {Twin} {Continuum}: A {Key} {Enabler}
for {Pervasive} {Cyber}-{Physical} {Environments},''} in \emph{2024 33rd
{International} {Conference} on {Computer} {Communications} and
{Networks} ({ICCCN})}, Jul. 2024, pp. 1--9. doi:
\href{https://doi.org/10.1109/ICCCN61486.2024.10637565}{10.1109/ICCCN61486.2024.10637565}.}

\bibitem[\citeproctext]{ref-alcaraz_digital_2022}
\CSLLeftMargin{{[}16{]} }%
\CSLRightInline{C. Alcaraz and J. Lopez, {``Digital {Twin}: {A}
{Comprehensive} {Survey} of {Security} {Threats},''} \emph{IEEE
Communications Surveys \& Tutorials}, vol. 24, no. 3, pp. 1475--1503,
2022, doi:
\href{https://doi.org/10.1109/COMST.2022.3171465}{10.1109/COMST.2022.3171465}.}

\bibitem[\citeproctext]{ref-kulik_security_2024}
\CSLLeftMargin{{[}17{]} }%
\CSLRightInline{T. Kulik, Z. Kazemi, and P. G. Larsen, {``Security and
{Privacy}-related {Issues} in a {Digital} {Twin} {Context},''} in
\emph{The {Engineering} of {Digital} {Twins}}, J. Fitzgerald, C. Gomes,
and P. G. Larsen, Eds., Cham: Springer International Publishing, 2024,
pp. 313--344. doi:
\href{https://doi.org/10.1007/978-3-031-66719-0_13}{10.1007/978-3-031-66719-0\_13}.}

\bibitem[\citeproctext]{ref-talasila_realising_2024}
\CSLLeftMargin{{[}18{]} }%
\CSLRightInline{P. Talasila, P. H. Mikkelsen, S. Gil, and P. G. Larsen,
{``Realising {Digital} {Twins},''} in \emph{The {Engineering} of
{Digital} {Twins}}, J. Fitzgerald, C. Gomes, and P. G. Larsen, Eds.,
Cham: Springer International Publishing, 2024, pp. 225--256. doi:
\href{https://doi.org/10.1007/978-3-031-66719-0_11}{10.1007/978-3-031-66719-0\_11}.}

\bibitem[\citeproctext]{ref-talasila_composable_2025}
\CSLLeftMargin{{[}19{]} }%
\CSLRightInline{P. Talasila \emph{et al.}, {``Composable digital twins
on {Digital} {Twin} as a {Service} platform,''} \emph{SIMULATION}, vol.
101, no. 3, pp. 287--311, Mar. 2025, doi:
\href{https://doi.org/10.1177/00375497241298653}{10.1177/00375497241298653}.}

\bibitem[\citeproctext]{ref-gil_survey_2024}
\CSLLeftMargin{{[}20{]} }%
\CSLRightInline{S. Gil, P. H. Mikkelsen, C. Gomes, and P. G. Larsen,
{``Survey on open‐source digital twin frameworks--{A} case study
approach,''} \emph{Software: Practice and Experience}, vol. 54, no. 6,
pp. 929--960, Jun. 2024, doi:
\href{https://doi.org/10.1002/spe.3305}{10.1002/spe.3305}.}

\bibitem[\citeproctext]{ref-zech_digital-twins-as-x2d-service_2024}
\CSLLeftMargin{{[}21{]} }%
\CSLRightInline{P. Zech, C. Nardin, S. Ristov, M. Flora, and R. Breu,
{``Digital-{Twins}-as-a-{Service} in {Construction} {Engineering},''} in
\emph{2024 {IEEE} 20th {International} {Conference} on {Automation}
{Science} and {Engineering} ({CASE})}, Aug. 2024, pp. 3004--3010. doi:
\href{https://doi.org/10.1109/CASE59546.2024.10711409}{10.1109/CASE59546.2024.10711409}.}

\bibitem[\citeproctext]{ref-duran_toward_2026}
\CSLLeftMargin{{[}22{]} }%
\CSLRightInline{K. Duran \emph{et al.}, {``Toward {Digital}
{Twin}-as-a-{Service} ({DTaaS}) {Platforms}: {A} {Survey} on
{Architecture}, {Design} {Requirements}, and {Performance} {Metrics},''}
\emph{IEEE Communications Surveys \& Tutorials}, vol. 28, pp.
1845--1878, 2026, doi:
\href{https://doi.org/10.1109/COMST.2025.3635582}{10.1109/COMST.2025.3635582}.}

\bibitem[\citeproctext]{ref-frasheri_addressing_2023}
\CSLLeftMargin{{[}23{]} }%
\CSLRightInline{M. Frasheri \emph{et al.}, {``Addressing time
discrepancy between digital and physical twins,''} \emph{Robotics and
Autonomous Systems}, vol. 161, p. 104347, Mar. 2023, doi:
\href{https://doi.org/10.1016/j.robot.2022.104347}{10.1016/j.robot.2022.104347}.}

\bibitem[\citeproctext]{ref-zhang_knowledge_2024}
\CSLLeftMargin{{[}24{]} }%
\CSLRightInline{N. Zhang, R. Bahsoon, N. Tziritas, and G.
Theodoropoulos, {``Knowledge {Equivalence} in {Digital} {Twins} of
{Intelligent} {Systems},''} \emph{ACM Transactions on Modeling and
Computer Simulation}, vol. 34, no. 1, pp. 1--37, Jan. 2024, doi:
\href{https://doi.org/10.1145/3635306}{10.1145/3635306}.}

\bibitem[\citeproctext]{ref-gomes_sensing_2024}
\CSLLeftMargin{{[}25{]} }%
\CSLRightInline{C. Gomes, D. E. L. Rötter, A. Iosifidis, H. Feng, H.
Ejersbo, and M. Frasheri, {``Sensing and {Communication} of {Data} from
the {Physical} {Twin},''} in \emph{The {Engineering} of {Digital}
{Twins}}, J. Fitzgerald, C. Gomes, and P. G. Larsen, Eds., Cham:
Springer International Publishing, 2024, pp. 147--171. doi:
\href{https://doi.org/10.1007/978-3-031-66719-0_7}{10.1007/978-3-031-66719-0\_7}.}

\bibitem[\citeproctext]{ref-oakes_case_2024}
\CSLLeftMargin{{[}26{]} }%
\CSLRightInline{B. J. Oakes \emph{et al.}, {``Case {Studies} in
{Digital} {Twins},''} in \emph{The {Engineering} of {Digital} {Twins}},
J. Fitzgerald, C. Gomes, and P. G. Larsen, Eds., Cham: Springer
International Publishing, 2024, pp. 257--310. doi:
\href{https://doi.org/10.1007/978-3-031-66719-0_12}{10.1007/978-3-031-66719-0\_12}.}

\bibitem[\citeproctext]{ref-committee_on_foundational_research_gaps_and_future_directions_for_digital_twins_foundational_2024}
\CSLLeftMargin{{[}27{]} }%
\CSLRightInline{\emph{Foundational {Research} {Gaps} and {Future}
{Directions} for {Digital} {Twins}}. Washington, D.C.: National
Academies Press, 2024. doi:
\href{https://doi.org/10.17226/26894}{10.17226/26894}.}

\bibitem[\citeproctext]{ref-hugues_twinops_2022}
\CSLLeftMargin{{[}28{]} }%
\CSLRightInline{J. Hugues, J. Yankel, J. Hudak, and A. Hristozov,
{``Twinops: {Digital} twins meets devops,''} \emph{CARNEGIE-MELLON UNIV
PITTSBURGH PA, Tech. Rep.}, 2022, Accessed: Jan. 08, 2025. {[}Online{]}.
Available: \url{https://apps.dtic.mil/sti/trecms/pdf/AD1168471.pdf}}

\bibitem[\citeproctext]{ref-barbie_toward_2024}
\CSLLeftMargin{{[}29{]} }%
\CSLRightInline{A. Barbie and W. Hasselbring, {``Toward
{Reproducibility} of {Digital} {Twin} {Research}: {Exemplified} with the
{PiCar}-{X}.''} arXiv, Aug. 2024. doi:
\href{https://doi.org/10.48550/arXiv.2408.13866}{10.48550/arXiv.2408.13866}.}

\bibitem[\citeproctext]{ref-dalibor_cross-domain_2022}
\CSLLeftMargin{{[}30{]} }%
\CSLRightInline{M. Dalibor \emph{et al.}, {``A {Cross}-{Domain}
{Systematic} {Mapping} {Study} on {Software} {Engineering} for {Digital}
{Twins},''} \emph{Journal of Systems and Software}, vol. 193, p. 111361,
Nov. 2022, doi:
\href{https://doi.org/10.1016/j.jss.2022.111361}{10.1016/j.jss.2022.111361}.}

\bibitem[\citeproctext]{ref-mertens_continuous_2024}
\CSLLeftMargin{{[}31{]} }%
\CSLRightInline{J. Mertens, S. Klikovits, F. Bordeleau, J. Denil, and Ø.
Haugen, {``Continuous {Evolution} of {Digital} {Twins} using the
{DarTwin} {Notation}.''} arXiv, Oct. 2024. Accessed: Nov. 11, 2024.
{[}Online{]}. Available: \url{http://arxiv.org/abs/2410.23389}}

\bibitem[\citeproctext]{ref-beaumont_towards_2025}
\CSLLeftMargin{{[}32{]} }%
\CSLRightInline{G. Beaumont, A. Beugnard, S. Martínez, C. Urtado, and S.
Vauttier, {``Towards {Automating} the {Life} {Cycle} {Management} of
{Digital} {Twins},''} in \emph{{ER2025}-44th {International}
{Conference} on {Conceptual} {Modeling}}, 2025. Accessed: Oct. 22, 2025.
{[}Online{]}. Available: \url{https://hal.science/hal-05225783/}}

\bibitem[\citeproctext]{ref-kamburjan_declarative_2024}
\CSLLeftMargin{{[}33{]} }%
\CSLRightInline{E. Kamburjan, N. Bencomo, S. L. Tapia Tarifa, and E. B.
Johnsen, {``Declarative {Lifecycle} {Management} in {Digital}
{Twins},''} in \emph{Proceedings of the {ACM}/{IEEE} 27th
{International} {Conference} on {Model} {Driven} {Engineering}
{Languages} and {Systems}}, Linz Austria: ACM, Sep. 2024, pp. 353--363.
doi:
\href{https://doi.org/10.1145/3652620.3688248}{10.1145/3652620.3688248}.}

\bibitem[\citeproctext]{ref-kamburjan_greenhousedt_2024}
\CSLLeftMargin{{[}34{]} }%
\CSLRightInline{E. Kamburjan \emph{et al.}, {``{GreenhouseDT}: {An}
{Exemplar} for {Digital} {Twins},''} in \emph{Proceedings of the 19th
{International} {Symposium} on {Software} {Engineering} for {Adaptive}
and {Self}-{Managing} {Systems}}, in {SEAMS} '24. New York, NY, USA:
Association for Computing Machinery, Jun. 2024, pp. 175--181. doi:
\href{https://doi.org/10.1145/3643915.3644108}{10.1145/3643915.3644108}.}

\bibitem[\citeproctext]{ref-goffi_engineering_2025}
\CSLLeftMargin{{[}35{]} }%
\CSLRightInline{P.-E. Goffi, R. Tremblay, and B. Oakes, {``Engineering a
{Digital} {Twin} for the {Monitoring} and {Control} of {Beer}
{Fermentation} {Sampling}.''} arXiv, Aug. 2025. doi:
\href{https://doi.org/10.48550/arXiv.2508.18452}{10.48550/arXiv.2508.18452}.}

\bibitem[\citeproctext]{ref-gomes_digital_2025}
\CSLLeftMargin{{[}36{]} }%
\CSLRightInline{C. Gomes \emph{et al.}, {``Digital {Twin} {Tutorial}:
{The} {Incubator} {Case} {Study},''} in \emph{Engineering {Trustworthy}
{Software} {Systems}: 6th {International} {School}, {SETSS} 2024,
{Chongqing}, {China}, {April} 14--21, 2024, {Tutorial} {Lectures}}, J.
P. Bowen, C. Gomes, and Z. Liu, Eds., Singapore: Springer Nature, 2025,
pp. 68--101. doi:
\href{https://doi.org/10.1007/978-981-96-4656-2_3}{10.1007/978-981-96-4656-2\_3}.}

\end{CSLReferences}

\end{document}
